\documentclass[11pt]{article}

\usepackage[T1]{fontenc}
\usepackage[utf8]{inputenc}
\usepackage{lmodern}
\usepackage[margin=1in]{geometry}
\usepackage{microtype}
\usepackage{amsmath,amssymb,mathtools,bm}
\usepackage{booktabs,tabularx,array}
\usepackage{enumitem}
\usepackage{xcolor}
\usepackage{hyperref}

\hypersetup{
  colorlinks=true,
  linkcolor=blue!55!black,
  urlcolor=blue!55!black,
  pdftitle={BCC nl Eshelby Hybrid Mobility Law}
}

\setcounter{tocdepth}{2}
\setlist{nosep,leftmargin=*}
\allowdisplaybreaks

\newcolumntype{Y}{>{\raggedright\arraybackslash}X}
\newcolumntype{P}[1]{>{\raggedright\arraybackslash}p{#1}}
\newcommand{\vect}[1]{\bm{#1}}
\newcommand{\norm}[1]{\left\lVert #1\right\rVert}
\newcommand{\abs}[1]{\left\lvert #1\right\rvert}
\newcommand{\dyad}[2]{#1\!\otimes\!#2}
\newcommand{\code}[1]{\texttt{\detokenize{#1}}}
\newcommand{\dd}{\,\mathrm{d}}

\title{The \texttt{BCC\_nl\_Eshelby} Hybrid Mobility Law\\
       Source Guide and Mathematical Formulation}
\author{KGParadis source-level technical documentation}
\date{Code snapshot reviewed 7 August 2026}

\begin{document}

\maketitle

\begin{abstract}
The mobility selector \code{BCC_nl_Eshelby} combines two existing BCC
constitutive laws on each arm attached to a node.  Portions outside the union
of Eshelby inclusions use the nonlinear screw and linear edge/junction model
from \path{MobilityLaw_BCC_nl.cc}; portions inside use exactly the
character-dependent drag tensor from \path{MobilityLaw_BCC_0.cc}.  This
document derives the inclusion weight, both constitutive branches, their
hybrid residual and Jacobian, and the corresponding source and build wiring.
The implementation is a complete, independent translation unit; the original
\path{MobilityLaw_BCC_nl.cc} is not modified.
\end{abstract}

\tableofcontents
\newpage

\section{Purpose and source layout}

The public entry points are
\begin{align*}
 &\code{InitMob_BCC_nl_Eshelby(Home_t *)},\\
 &\code{Mobility_BCC_nl_Eshelby(Home_t *,Node_t *,MobArgs_t *)}.
\end{align*}
They are implemented in
\path{src/MobilityLaw_BCC_nl_Eshelby.cc}.  The control-file selector is
\begin{center}
\code{mobilityLaw = "BCC_nl_Eshelby"}.
\end{center}
The source is compiled only when \code{ESHELBY} is defined.  It contains its
own copies of the private BCC nonlinear helpers and its own file-static
constants \code{Linear}, \code{NonLinear}, and \code{V0}.  Its initializer
must therefore compute those constants itself; calling \code{InitMob_BCC_nl}
would initialize a different translation unit's private state.

The implementation and registration are divided as follows.

\begin{center}
\small
\begin{tabularx}{\textwidth}{@{}P{0.39\textwidth}Y@{}}
\toprule
File & Responsibility \\
\midrule
\path{src/MobilityLaw_BCC_nl_Eshelby.cc} &
  Complete mobility implementation, inclusion interval union, BCC\_0 tensor,
  hybrid Newton residual and Jacobian. \\
\path{include/Mobility.h} &
  Guarded enumeration, prototypes, and selector attributes. \\
\path{src/DisableUnneededParams.cc} &
  Retains the union of BCC\_nl, BCC\_0, pressure, and inclusion parameters in
  restart/control output. \\
\path{makefile.srcs} and \path{src/makefile.dep} &
  Object registration and include dependencies. \\
\bottomrule
\end{tabularx}
\end{center}

\section{Arm geometry and notation}

For an arm from the current node at \(\vect{x}_1\) to a neighboring node,
the minimum-image segment vector and its length are
\begin{equation}
 \vect{r}=\operatorname{ZImage}(\vect{x}_2-\vect{x}_1),\qquad
 L=\norm{\vect{r}},\qquad \vect{t}=\frac{\vect{r}}{L}.
 \label{eq:arm}
\end{equation}
Let \(\vect{b}\) be the arm Burgers vector.  The BCC nonlinear decomposition
uses the projected screw and edge lengths
\begin{equation}
 L_s=\frac{\abs{\vect{r}\mathbin{\cdot}\vect{b}}}{\norm{\vect{b}}},
 \qquad
 L_e=\sqrt{\max(0,L^2-L_s^2)}.
 \label{eq:lengths}
\end{equation}
Every arm contribution has the usual nodal factor \(1/2\).  Vectors and
tensors below are evaluated in the crystal frame used by the copied nonlinear
solver; Section~\ref{sec:frames} describes frame conversion.

The source classifies an arm as a product-zero or junction Burgers type when
\begin{equation}
 \abs{b_x b_y b_z}<10^{-12}.
 \label{eq:junction-test}
\end{equation}
Other Burgers vectors use the regular BCC \(\langle111\rangle\) branch.

\section{Eshelby coverage of an arm}

\subsection{Ellipsoid intersection and periodic images}

For an inclusion centered at \(\vect{C}\), with semi-axes
\(a_1,a_2,a_3\) and orthonormal principal directions
\(\vect{q}_1,\vect{q}_2,\vect{q}_3\), its interior is
\begin{equation}
 \sum_{k=1}^{3}
 \left(\frac{\vect{q}_k\mathbin{\cdot}(\vect{x}-\vect{C})}{a_k}\right)^2
 \leq 1.
 \label{eq:ellipsoid}
\end{equation}
The code first moves \(\vect{C}\) to the periodic image nearest
\(\vect{x}_1\).  It then intersects
\begin{equation}
 \vect{x}(s)=\vect{x}_1+s\vect{r},\qquad 0\leq s\leq1,
 \label{eq:parametric-arm}
\end{equation}
with the ellipsoid by transforming it to a unit sphere.  Each intersected
particle produces a clamped interval \([a_k,b_k]\subseteq[0,1]\).  Tangent or
zero-width intervals are discarded.  The intervals are sorted and merged so
overlapping inclusions are counted only once.  Denote their union by
\(U\).

\subsection{Nodal rather than geometric length fraction}

The inside weight is not simply the covered length.  For the current endpoint
the linear nodal shape function is \(N_1(s)=1-s\), so the normalized covered
weight is
\begin{equation}
 \boxed{
 \phi=2\int_U(1-s)\dd s
 =\sum_{[a_k,b_k]\subset U}
   \left[(1-a_k)^2-(1-b_k)^2\right] .}
 \label{eq:phi}
\end{equation}
Thus \(0\leq\phi\leq1\), with \(\phi=0\) for an uncovered arm and
\(\phi=1\) for a wholly covered arm.  Reversing the arm at its neighboring
endpoint naturally changes the shape-function weight.  This is why the law
does not perform a binary switch based only on whether the node position is
inside a particle.

If \code{inclusionFile} is empty, \code{enableInclusions} is false and the
helper returns \(\phi=0\) without doing intersection work.

\section{Outside law: original BCC nonlinear mobility}

\subsection{Temperature- and pressure-dependent screw constants}

Write
\[
 p=\code{MobExpP},\quad q=\code{MobExpQ},\quad
 \alpha=\code{MobAlpha},\quad
 H=\frac{\code{MobDeltaH0}}{k_B T},
\]
and let \(P_{\mathrm{GPa}}=P/10^9\).  The initializer computes
\begin{align}
 \mathcal{L}
 &=A C_0\left(1-C\frac{T}{T_m}+D P_{\mathrm{GPa}}\right),
 \label{eq:linear-constant}\\
 y&=\alpha^p,\nonumber\\
 \mathcal{N}
 &=\frac{\mathcal{L}}
 {\exp[-H(1-y)^q]\,[1+pqH(1-y)^{q-1}y]},
 \label{eq:nonlinear-constant}\\
 V_0&=\alpha\mathcal{L}
      -\alpha\mathcal{N}\exp[-H(1-y)^q].
 \label{eq:v0-constant}
\end{align}
Here \(A,C_0,C,D,T_m\) correspond to \code{MobCoeffA},
\code{MobCoeffC0}, \code{MobCoeffC}, \code{MobCoeffD}, and
\code{meltTemp}.  The per-Burgers-magnitude quantities used by the force law
are
\begin{equation}
 M_L=\frac{\mathcal{L}}{b_{\rm mag}},\qquad
 M_N=\frac{\mathcal{N}}{b_{\rm mag}},\qquad
 v_0=\frac{V_0}{b_{\rm mag}}.
 \label{eq:scaled-constants}
\end{equation}

For one screw Burgers family, define
\begin{equation}
 f_P=\code{MobPeierls}\,\norm{\vect{b}},\qquad
 \omega=\frac{M_N}{f_P},\qquad
 v_{\rm tr}=\alpha M_L-v_0.
\end{equation}
For \(0<v<v_{\rm tr}\), the scalar glide force \(f(v)\) is the numerical
inverse of
\begin{equation}
 \boxed{v=\omega f
 \exp\!\left[-H\left(1-(f/f_P)^p\right)^q\right].}
 \label{eq:thermal-law}
\end{equation}
For \(v\geq v_{\rm tr}\), the continuation is linear,
\begin{equation}
 f(v)=\frac{(v+v_0)f_P}{M_L}.
 \label{eq:high-speed-law}
\end{equation}
The source solves Equation~\eqref{eq:thermal-law} with its inherited Newton
iterations and also returns \(f'(v)\).

\subsection{Screw force and linear edge/junction tensors}

Let \(\widehat{\vect{b}}=\vect{b}/\norm{\vect{b}}\),
\(P_b=I-\dyad{\widehat{\vect{b}}}{\widehat{\vect{b}}}\), and
\begin{equation}
 \vect{v}_{\perp}=P_b\vect{v},\qquad
 v_{\perp}=\norm{\vect{v}_{\perp}},\qquad
 \widehat{\vect{v}}_{\perp}=\vect{v}_{\perp}/v_{\perp}.
\end{equation}
The nonlinear screw force per unit screw length is
\begin{equation}
 \vect{F}_s(\vect{v};\vect{b})
 =B_l\dyad{\widehat{\vect{b}}}{\widehat{\vect{b}}}\vect{v}
  +f(v_{\perp})\widehat{\vect{v}}_{\perp},
 \qquad B_l=\frac{1}{\code{MobLine}}.
 \label{eq:screw-force}
\end{equation}
Away from \(v_{\perp}=0\), its derivative is
\begin{equation}
 D_s=B_l\dyad{\widehat{\vect{b}}}{\widehat{\vect{b}}}
 +f'(v_{\perp})
   \dyad{\widehat{\vect{v}}_{\perp}}{\widehat{\vect{v}}_{\perp}}
 +\frac{f(v_{\perp})}{v_{\perp}}
  \left(P_b-
   \dyad{\widehat{\vect{v}}_{\perp}}{\widehat{\vect{v}}_{\perp}}
  \right).
 \label{eq:screw-jacobian}
\end{equation}
The implementation uses the finite zero-velocity limit of this expression.

For a regular arm, set
\(\vect{n}=(\vect{b}\times\vect{t})/
\norm{\vect{b}\times\vect{t}}\) and use the orthogonal line direction
\(\vect{t}\) up to an irrelevant sign.  Its linear edge tensor is
\begin{equation}
 D_e=B_g\dyad{\widehat{\vect{b}}}{\widehat{\vect{b}}}
     +B_c\dyad{\vect{n}}{\vect{n}}
     +B_l\dyad{\vect{t}}{\vect{t}},
 \quad
 B_g=\frac1{\code{MobGlide}},\quad
 B_c=\frac1{\code{MobClimb}}.
 \label{eq:edge-tensor}
\end{equation}
For a product-zero/junction arm, the whole arm uses
\begin{equation}
 D_j=B_c(I-\dyad{\vect{t}}{\vect{t}})
     +B_l\dyad{\vect{t}}{\vect{t}}.
 \label{eq:junction-tensor}
\end{equation}

\section{Inside law: exact BCC\_0 arm tensor}

The inside portion does not use a scalar Eshelby resistance multiplier.  It
uses the character-dependent tensor from \path{MobilityLaw_BCC_0.cc}.  Define
\begin{equation}
 B_s=\frac1{\code{MobScrew}},\qquad
 B_e=\frac1{\code{MobEdge}},\qquad
 B_{ec}=\frac1{\code{MobClimb}},\qquad
 B_{l0}=10^{-2}\min(B_s,B_e),
 \label{eq:bcc0-coeffs}
\end{equation}
and
\begin{equation}
 c^2=\frac{(\vect{t}\mathbin{\cdot}\vect{b})^2}{\norm{\vect{b}}^2}.
 \label{eq:character}
\end{equation}

For a product-zero Burgers vector, Equation~\eqref{eq:junction-test}, the
BCC\_0 tensor is
\begin{equation}
 \boxed{B^0=B_{ec}I+(B_{l0}-B_{ec})
 \dyad{\vect{t}}{\vect{t}}.}
 \label{eq:bcc0-special}
\end{equation}
For a regular near-screw arm, \(1-c^2\leq10^{-12}\), it is
\begin{equation}
 \boxed{B^0=B_sI+(B_{l0}-B_s)\dyad{\vect{t}}{\vect{t}}.}
 \label{eq:bcc0-screw}
\end{equation}

For a mixed regular arm, form the orthonormal directions
\begin{equation}
 \vect{n}=\frac{\vect{b}\times\vect{t}}
 {\sqrt{(1-c^2)\norm{\vect{b}}^2}},
 \qquad \vect{m}=\vect{n}\times\vect{t},
 \label{eq:bcc0-basis}
\end{equation}
and the interpolated glide and climb drags
\begin{align}
 B_g^0&=\left[B_e^{-2}+(B_s^{-2}-B_e^{-2})c^2\right]^{-1/2},
 \label{eq:bcc0-glide}\\
 B_c^0&=\left[B_{ec}^{2}+(B_s^{2}-B_{ec}^{2})c^2\right]^{1/2}.
 \label{eq:bcc0-climb}
\end{align}
The tensor assembled by the source is equivalently
\begin{equation}
 \boxed{B^0=B_{l0}\dyad{\vect{t}}{\vect{t}}
             +B_c^0\dyad{\vect{n}}{\vect{n}}
             +B_g^0\dyad{\vect{m}}{\vect{m}}.}
 \label{eq:bcc0-mixed}
\end{equation}
At pure edge character, \(c^2=0\), this reduces to line, climb, and edge
drags \(B_{l0},B_{ec},B_e\) along \(\vect{t},\vect{n},\vect{m}\).

\section{Hybrid residual and Jacobian}

For a regular arm \(i\), the drag-force contribution assigned to the current
node is
\begin{equation}
 \boxed{
 \vect{R}_i^{\rm drag}=
 \frac12(1-\phi_i)
 \left[L_{e,i}D_{e,i}\vect{v}
       +L_{s,i}\vect{F}_s(\vect{v};\vect{b}_i)\right]
 +\frac12\phi_iL_iB_i^0\vect{v}.}
 \label{eq:hybrid-regular}
\end{equation}
For a product-zero/junction arm it is
\begin{equation}
 \boxed{
 \vect{R}_i^{\rm drag}=
 \frac12L_i\left[(1-\phi_i)D_{j,i}+\phi_iB_i^0\right]\vect{v}.}
 \label{eq:hybrid-junction}
\end{equation}
Notice that the BCC\_0 part owns the full physical arm length \(L_i\); it is
not decomposed into \(L_s\) and \(L_e\).

Collinear screw Burgers vectors are grouped.  Define
\begin{equation}
 C_g=\frac12\sum_{i\in g}(1-\phi_i)L_{s,i}
 \label{eq:screw-group}
\end{equation}
and let \(D_{\rm lin}\) contain all edge/junction terms plus all inside
BCC\_0 tensors.  The Newton solve uses the force residual
\begin{equation}
 \boxed{
 \vect{r}(\vect{v})=\vect{F}_{\rm node}
 -D_{\rm lin}\vect{v}
 -\sum_g C_g\vect{F}_s(\vect{v};\vect{b}_g)=\vect{0},}
 \label{eq:residual}
\end{equation}
with Jacobian
\begin{equation}
 \boxed{
 J(\vect{v})=-D_{\rm lin}-\sum_g C_gD_s(\vect{v};\vect{b}_g).}
 \label{eq:jacobian}
\end{equation}
The iteration applies \(\vect{v}\leftarrow\vect{v}-J^{-1}\vect{r}\), using
the convergence and correction controls inherited from BCC\_nl.

The original BCC\_nl one-screw shortcut solves only a two-dimensional
nonlinear subspace and restores the Burgers-direction velocity analytically.
A mixed BCC\_0 tensor need not have \(\vect{b}\) as an eigenvector.  Therefore
the new implementation disables that shortcut whenever any arm has
\(\phi_i>0\), and uses the full three-dimensional Newton system instead.

\section{Frames, outputs, and limiting cases}
\label{sec:frames}

Intersection geometry is evaluated in the simulation/laboratory frame:
the minimum-image arm \(\vect{r}_{\rm lab}\), node position, and nearest
periodic inclusion center are passed to the ellipsoid routine.  A saved copy
of \(\vect{r}_{\rm lab}\) isolates this calculation from crystallographic
rotation.

When \code{useLabFrame} is enabled, force, velocity, arm direction, and Burgers
vector are rotated with \code{rotMatrixInverse}; all drag tensors and the
Newton solve are then evaluated in the crystal frame.  The converged velocity
is rotated back with \code{rotMatrix}.  For an orthogonal rotation \(Q\), this
is the covariance relation
\begin{equation}
 B^0_{\rm crystal}=Q^TB^0_{\rm lab}Q.
\end{equation}

The important limits are:
\begin{align}
 \phi_i=0\ \forall i
 &\quad\Longrightarrow\quad
 \text{the copied BCC\_nl algebra and execution path},
 \label{eq:limit-out}\\
 \phi_i=1\ \forall i
 &\quad\Longrightarrow\quad
 D_{\rm lin}=\sum_i\frac12L_iB_i^0,\qquad C_g=0.
 \label{eq:limit-in}
\end{align}
Thus complete coverage gives the same nodal velocity as
\code{Mobility_BCC_0} when \code{includeInertia=0}.  The nonlinear parent law
does not contain BCC\_0's optional mass term, so inertial parity is not claimed
when \code{includeInertia=1}.

The \code{mobArgs->invDragMatrix} output intentionally follows the inherited
BCC\_nl convention: it is an inverse of the negative Newton Jacobian in the
solver frame.  Even at complete coverage, it is not the positive lab-frame
\(B^{-1}\) stored by \code{Mobility_BCC_0}; the nodal velocity is nevertheless
the BCC\_0 velocity under the non-inertial condition above.

\section{Parameters and control-file use}

\begin{center}
\small
\begin{tabularx}{\textwidth}{@{}P{0.29\textwidth}P{0.25\textwidth}Y@{}}
\toprule
Parameter group & Used by & Meaning \\
\midrule
\code{inclusionFile} & Geometry & Inclusion centers, radii, orientations, and
candidate intersection data. \\
\code{MobScrew}, \code{MobEdge}, \code{MobClimb} & BCC\_0 inside & Mobilities
whose reciprocals define Equation~\eqref{eq:bcc0-coeffs}. \\
\code{MobGlide}, \code{MobLine}, \code{MobClimb} & BCC\_nl outside & Linear
edge, line, and climb drags. \\
\code{MobPeierls}, \code{MobExpP}, \code{MobExpQ}, \code{MobAlpha},
\code{MobDeltaH0} & BCC\_nl outside & Thermally activated screw law. \\
\code{MobCoeffA}, \code{MobCoeffC0}, \code{MobCoeffC}, \code{MobCoeffD},
\code{meltTemp}, \code{TempK}, \code{pressure} & BCC\_nl outside & Initializer
constants in Equations~\eqref{eq:linear-constant}--\eqref{eq:v0-constant}. \\
\code{MobFileNL}, \code{ecFile}, \code{meltTempFile} & Existing nonlinear
initialization paths & Optional table/file inputs retained consistently with
BCC\_nl. \\
\code{MobEshelbyResist} & Not used & Deliberately disabled for this selector;
the inside branch is the literal BCC\_0 tensor, not a scalar-rescaled outside
law. \\
\bottomrule
\end{tabularx}
\end{center}

The following is a ready-to-use control-file example based on
\path{tests/screwEsh.ctrl}.  It includes the simulation controls needed by that
test as well as the complete parameter union used by the hybrid law.  The
material constants shown are the existing tantalum defaults in
\path{src/Param.cc}; they should be replaced by the calibration for the
material being simulated.

\subsection{Complete example control input}

\begin{verbatim}
dirname = "tests/screwEsh_nl_Eshelby_results"
inclusionFile = "tests/screwEsh.dat"

# Domain and cell geometry
numXdoms = 1
numYdoms = 1
numZdoms = 1
numXcells = 3
numYcells = 3
numZcells = 3

# Discretization and topological changes
maxSeg = 50
minSeg = 15
remeshRule = 2
enforceGlidePlanes = 0
enableCrossSlip = 0
splitMultiNodeFreq = 5

# Time integration
maxstep = 1000000
timestepIntegrator = "trapezoid"
nextDT = 1.0e-12
rTol = 0.5

# Fast multipole controls
fmEnabled = 0
fmMPOrder = 2
fmTaylorOrder = 5

# Loading and thermodynamic state
loadType = 1
eRate = 1.0e5
indxErate = 2
edotdir = [
  0.0
  0.0
  1.0
]
TempK = 600.0
pressure = 0.0

# Select the hybrid law
mobilityLaw = "BCC_nl_Eshelby"
materialTypeName = "BCC"

# Elastic and core parameters used by the simulation
shearModulus = 8.600000e10
pois = 2.910000e-1
burgMag = 2.480000e-10
YoungModulus = 2.220520e11
rc = 2.0
rann = 5.0
collisionMethod = 3

# BCC_0 mobility used on inclusion-covered arm portions
MobScrew = 1.0e2
MobEdge  = 1.0e2
MobClimb = 1.0e-5

# BCC_nl mobility used on uncovered arm portions
MobGlide = 1.0e2
MobLine = 1.0e4
MobPeierls = 3.240000e8
MobDeltaH0 = 1.728748e-19
MobAlpha = 9.661236e-1
MobExpP = 0.5
MobExpQ = 1.23
MobCoeffA = 1.525
MobCoeffC = 2.032370e-1
MobCoeffC0 = 2.048000e3
MobCoeffD = 1.626950e-4
meltTemp = 3.170683e3

# Optional nonlinear-property files; empty uses values above
MobFileNL = ""
ecFile = ""
meltTempFile = ""

# Output controls
fluxfile = 1
fluxfreq = 200
savecn = 1
savecnfreq = 200
saveprop = 1
savepropfreq = 200
savetimers = 0
velfile = 0
velfilefreq = 200
\end{verbatim}

The associated nodal data and particle file for this example are
\path{tests/screwEsh.data} and \path{tests/screwEsh.dat}.  A one-step serial
smoke run can be launched as
\begin{verbatim}
bin/paradis -maxstep 1 -s -d tests/screwEsh.data \
  path/to/BCC_nl_Eshelby.ctrl
\end{verbatim}

\subsection{Inputs specific to the hybrid switch}

Only three changes are required when converting an existing BCC\_nl control
file:
\begin{enumerate}
 \item set \code{mobilityLaw = "BCC_nl_Eshelby"};
 \item provide a nonempty \code{inclusionFile};
 \item provide positive \code{MobScrew} and \code{MobEdge} values for the
       BCC\_0 inside tensor.  \code{MobClimb} is already required by BCC\_nl
       and is shared by both branches.
\end{enumerate}
\code{MobEshelbyResist} is not used by this selector and should be omitted.

An empty \code{inclusionFile} is permitted; it makes the law reduce to
BCC\_nl everywhere.

\section{Build and verification}

The selector, prototypes, and implementation are guarded by \code{ESHELBY}.
The object remains in the common source list, like the other Eshelby mobility
objects; without \code{ESHELBY} it compiles to an empty translation unit and
the selector is absent.  No changes are required in \path{Param.cc},
\path{Param.h}, or \path{Initialize.cc}, because mobility selection and
function-pointer initialization are generic.

The implementation was checked in the following ways:
\begin{itemize}
 \item full serial builds with and without \code{ESHELBY};
 \item a one-step ParaDiS run using \code{BCC_nl_Eshelby} and an inclusion
       file;
 \item symbolic comparison of every BCC\_0 tensor branch with
       \path{MobilityLaw_BCC_0.cc};
 \item \(10^5\) randomized orientations over wide mobility ranges, with only
       floating-point ordering differences (approximately \(10^{-12}\) at
       deliberately extreme ranges and near machine precision ordinarily);
 \item direct checks of overlapping, tangent, disjoint, and fully contained
       inclusion intervals;
 \item algebraic checks of the \(\phi=0\) and \(\phi=1\) limits.
\end{itemize}

\section{Implementation cautions}

\begin{enumerate}
 \item ``Inside'' is arm- and endpoint-weighted through Equation~\eqref{eq:phi};
       a node center inside an inclusion does not imply that all its arms have
       \(\phi=1\).
 \item The three BCC\_0 input mobilities must be finite and positive because
       the source uses their reciprocals.
 \item Complete BCC\_0 equivalence assumes \code{includeInertia=0}.
 \item \code{MobEshelbyResist} has no effect on this law by design.
 \item The outside-only branch deliberately retains the numerical behavior of
       the original BCC\_nl implementation.  Baseline solver changes should be
       made separately in both laws if that behavior is later revised.
\end{enumerate}

\end{document}
